% GENERATED FILE. Edit canonical lessons, then run npm run generate:books.
% canonical-source-hash: fnv1a64:06968a574daf7bb3
% canonical-lessons: ES-C763-escultura, ES-C763-retrato, ES-C763-lienzo, ES-C763-pintor, ES-C763-fotografiar

\chapter{En el museo}
\label{ch:es-en-el-museo}
\hlchaptermodality{763}

\begin{chapteropening}
\textbf{By the end of this chapter:} \emph{I can talk about sculptures, portraits, canvases and painters, and about taking photos.}

Complete the last lesson of chapter 763: I can talk about sculptures, portraits, canvases and painters, and about taking photos.
\end{chapteropening}

\section[la escultura]{la escultura — the sculpture}

\label{lesson:ES-C763-escultura}



\begin{quote}
Before the new one: say the Spanish for the cover, then the Spanish for to leaf through.
\end{quote}

\subsection*{You'll want to know: la escultura}
\textbf{la escultura} — \textquotedblleft{}the sculpture\textquotedblright{}.

\begin{quote}
En el museo hay una escultura enorme.
\end{quote}

\begin{quote}
\emph{There is a huge sculpture in the museum.}
\end{quote}

\emph{El escultor} — the sculptor.

\begin{cousinweb}
\emph{escultura} comes from Latin \emph{sculptūra} \textquotedblleft{}sculpture\textquotedblright{} (EtymDB), from \emph{sculpere} \textquotedblleft{}to carve in stone\textquotedblright{}.
\end{cousinweb}

\begin{grammarlens}[title={What you've built}]
One more word for telling, explaining and describing at B1.
\end{grammarlens}

\subsection*{Your turn}
\begin{itemize}
\raggedright
  \item \emph{Say it:} \emph{la escultura}
  \item \emph{Say it:} \emph{la escultura}, once more
  \item \emph{Say it:} a sentence of your own with \emph{la escultura}
  \item \emph{Recall:} say the Spanish for the cover, then the Spanish for to leaf through, then say \emph{la escultura} again
\end{itemize}

\begin{tcolorbox}[breakable,colback=teal!4,colframe=teal!35!black,title={Before you move on}]
What does la escultura mean? (\textquotedblleft{}The sculpture\textquotedblright{}.) Say it once more.
\end{tcolorbox}
\section[el retrato]{el retrato — the portrait}

\label{lesson:ES-C763-retrato}



\begin{quote}
Before the new one: say the Spanish for to leaf through, then the Spanish for the sculpture.
\end{quote}

\subsection*{You'll want to know: el retrato}
\textbf{el retrato} — \textquotedblleft{}the portrait\textquotedblright{}.

\begin{quote}
Este retrato es de mi abuela.
\end{quote}

\begin{quote}
\emph{This portrait is of my grandmother.}
\end{quote}

\emph{Un retrato robot} — an identikit picture.

\begin{cousinweb}
\emph{retrato} was borrowed from Italian \emph{ritratto} \textquotedblleft{}portrait\textquotedblright{} (EtymDB), from Latin \emph{retractus}, built on \emph{trahere} \textquotedblleft{}to draw, to pull\textquotedblright{}.
\end{cousinweb}

\begin{grammarlens}[title={What you've built}]
One more word for telling, explaining and describing at B1.
\end{grammarlens}

\subsection*{Your turn}
\begin{itemize}
\raggedright
  \item \emph{Say it:} \emph{el retrato}
  \item \emph{Say it:} \emph{el retrato}, once more
  \item \emph{Say it:} a sentence of your own with \emph{el retrato}
  \item \emph{Recall:} say the Spanish for to leaf through, then the Spanish for the sculpture, then say \emph{el retrato} again
\end{itemize}

\begin{tcolorbox}[breakable,colback=teal!4,colframe=teal!35!black,title={Before you move on}]
What does el retrato mean? (\textquotedblleft{}The portrait\textquotedblright{}.) Say it once more.
\end{tcolorbox}
\section[el lienzo]{el lienzo — the canvas}

\label{lesson:ES-C763-lienzo}



\begin{quote}
Before the new one: say the Spanish for the sculpture, then the Spanish for the portrait.
\end{quote}

\subsection*{You'll want to know: el lienzo}
\textbf{el lienzo} — \textquotedblleft{}the canvas\textquotedblright{}.

\begin{quote}
Compré un lienzo para pintar.
\end{quote}

\begin{quote}
\emph{I bought a canvas to paint on.}
\end{quote}

Also an old word for linen cloth.

\begin{cousinweb}
\emph{lienzo} is inherited from Latin \emph{linteum} \textquotedblleft{}linen cloth\textquotedblright{} (EtymDB), through Old Spanish \emph{lienço}.
\end{cousinweb}

\begin{grammarlens}[title={What you've built}]
One more word for telling, explaining and describing at B1.
\end{grammarlens}

\subsection*{Your turn}
\begin{itemize}
\raggedright
  \item \emph{Say it:} \emph{el lienzo}
  \item \emph{Say it:} \emph{el lienzo}, once more
  \item \emph{Say it:} a sentence of your own with \emph{el lienzo}
  \item \emph{Recall:} say the Spanish for the sculpture, then the Spanish for the portrait, then say \emph{el lienzo} again
\end{itemize}

\begin{tcolorbox}[breakable,colback=teal!4,colframe=teal!35!black,title={Before you move on}]
What does el lienzo mean? (\textquotedblleft{}The canvas\textquotedblright{}.) Say it once more.
\end{tcolorbox}
\section[el pintor]{el pintor — the painter}

\label{lesson:ES-C763-pintor}



\begin{quote}
Before the new one: say the Spanish for the portrait, then the Spanish for the canvas.
\end{quote}

\subsection*{You'll want to know: el pintor}
\textbf{el pintor} — \textquotedblleft{}the painter\textquotedblright{}.

\begin{quote}
Picasso fue un pintor muy famoso.
\end{quote}

\begin{quote}
\emph{Picasso was a very famous painter.}
\end{quote}

\emph{La pintora} for a woman. Also a house painter.

\begin{cousinweb}
\emph{pintor} is inherited from spoken Latin \emph{pinctor}, Latin \emph{pictor} \textquotedblleft{}painter\textquotedblright{} (EtymDB), from \emph{pingere} \textquotedblleft{}to paint\textquotedblright{}.
\end{cousinweb}

\begin{grammarlens}[title={What you've built}]
One more word for telling, explaining and describing at B1.
\end{grammarlens}

\subsection*{Your turn}
\begin{itemize}
\raggedright
  \item \emph{Say it:} \emph{el pintor}
  \item \emph{Say it:} \emph{el pintor}, once more
  \item \emph{Say it:} a sentence of your own with \emph{el pintor}
  \item \emph{Recall:} say the Spanish for the portrait, then the Spanish for the canvas, then say \emph{el pintor} again
\end{itemize}

\begin{tcolorbox}[breakable,colback=teal!4,colframe=teal!35!black,title={Before you move on}]
What does el pintor mean? (\textquotedblleft{}The painter\textquotedblright{}.) Say it once more.
\end{tcolorbox}
\section[fotografiar]{fotografiar — to photograph}

\label{lesson:ES-C763-fotografiar}



\begin{quote}
Before the new one: say the Spanish for the canvas, then the Spanish for the painter.
\end{quote}

\subsection*{You'll want to know: fotografiar}
\textbf{fotografiar} — \textquotedblleft{}to photograph\textquotedblright{}.

\begin{quote}
Está prohibido fotografiar los cuadros.
\end{quote}

\begin{quote}
\emph{Photographing the paintings is not allowed.}
\end{quote}

Everyday speech prefers \emph{hacer fotos} (Spain) or \emph{tomar fotos} (Latin America). Forms: \emph{fotografío}.

\begin{cousinweb}
\emph{fotografiar} is built on \emph{fotografía}, made from Greek \emph{phōs}, \emph{phōtós} \textquotedblleft{}light\textquotedblright{} and \emph{gráphein} \textquotedblleft{}to write\textquotedblright{}: writing with light.
\end{cousinweb}

\begin{grammarlens}[title={What you've built}]
One more word for telling, explaining and describing at B1.
\end{grammarlens}

\subsection*{Your turn}
\begin{itemize}
\raggedright
  \item \emph{Say it:} \emph{fotografiar}
  \item \emph{Say it:} \emph{fotografiar}, once more
  \item \emph{Say it:} a sentence of your own with \emph{fotografiar}
  \item \emph{Recall:} say the Spanish for the canvas, then the Spanish for the painter, then say \emph{fotografiar} again
\end{itemize}

\begin{tcolorbox}[breakable,colback=teal!4,colframe=teal!35!black,title={Before you move on}]
What does fotografiar mean? (\textquotedblleft{}To photograph\textquotedblright{}.) Say it once more.
\end{tcolorbox}
